\section{Introduction}
\label{sec:intro}

The rapid proliferation of task-specific deep neural networks has spurred immense interest in techniques that consolidate diverse capabilities into a single unified model. Traditionally, multi-task learning or sequential fine-tuning has been employed to achieve this goal. However, these methods suffer from severe limitations, including prohibitive training costs, catastrophic forgetting of previous tasks, and the requirement of concurrent access to all training datasets. 

To bypass these obstacles, \emph{model merging} has emerged as a powerful, training-free paradigm \cite{langley00, mitchell80}. Model merging seeks to fuse multiple independently fine-tuned task experts into a single, cohesive network directly in the parameter space. The simplest and most widely used baseline is \emph{Task Arithmetic} \cite{ilharco2022editing}, which computes the parameter difference (task vector) between each expert and the pretrained base model, then blends them via a uniform static scaling factor. While computationally trivial and elegant, Task Arithmetic assumes that parameter interactions across tasks are uniform and independent, leading to severe representation interference and performance degradation when task vectors collide.

To mitigate this interference, state-of-the-art methods have introduced \emph{adaptive model merging}, where merging coefficients are optimized on a layer-by-layer basis. Representative works such as AdaMerging \cite{yang2023adamerging}, RegCalMerge \cite{regcalmerge}, and PolyMerge \cite{polymerge} optimize these layer-wise coefficients at test time via \emph{test-time adaptation} (TTA). Specifically, they employ unsupervised entropy minimization on unlabeled target data streams to guide the parameter search. 

Despite their empirical success, this test-time optimization paradigm has several fundamental, unresolved theoretical and practical weaknesses:
\begin{itemize}
    \item \textbf{Heuristic Surrogate Objective:} Unsupervised entropy minimization is an indirect surrogate for the true multi-task classification objective. Minimizing the entropy of model predictions on unlabeled target data is not mathematically guaranteed to minimize the joint multi-task supervised loss.
    \item \textbf{Transductive Overfitting (The Overfitting-Optimizer Paradox):} Adapting coefficients on small, unlabeled batches of test data often causes the optimizer to overfit to the local transductive distribution, leading to catastrophic generalization collapse on unseen target samples.
    \item \textbf{Sacrificial Task Bias:} Without explicit joint objective constraints, test-time entropy minimization is highly susceptible to local minima that heavily favor one task at the severe expense of others (e.g., maximizing SVHN performance by entirely destroying MNIST capabilities).
    \item \textbf{Prohibitive Computational Overhead:} These methods require hundreds of forward-backward optimization steps at deployment time, which completely defeats the "training-free" promise of model merging and restricts their utility in latency-sensitive applications.
\end{itemize}

In this work, we approach model merging from a first-principles mathematical perspective, rejecting test-time heuristics in favor of formal loss-landscape geometry. We present \textbf{Curvature-Aware Analytical Model Merging (ACM)}, a mathematically rigorous, training-free framework that computes the optimal merging coefficients directly in a single step. 

Our core insight is that the joint multi-task loss can be formulated as a quadratic minimization problem via a second-order Taylor expansion around each individual task expert's local minimum. While modeling the full $D$-dimensional Hessian curvature of a modern neural network is a prohibitive $O(D^2)$ space and $O(D^3)$ time operation, we prove that by projecting the massive parameter space onto the extremely low-dimensional $K$-dimensional subspace spanned by the task vectors, we can capture the \emph{full, non-diagonal, cross-parameter Hessian curvature} along the directions of task updates with **zero diagonal approximation**. Unlike Fisher Merging \cite{matena2022merging}, which assumes a diagonal approximation of the Hessian to remain tractable, ACM maintains the full cross-parameter second-order interactions within the task-vector subspace. This projection reduces the optimization complexity to a trivial $O(K^2)$ memory and $O(K^3)$ inversion cost (where $K$ is the number of tasks, e.g., $K \ll D$). Under this framework, we derive an exact, closed-form analytical solution $\Lambda^{l, *} = (A^l + \gamma I)^{-1} (b^l - d^l)$ for each layer $l$ under Ridge regularization, where $d^l$ is the linear projection of expert gradients.

To compute the elements of the projected Hessian matrix $A^l$ and vector $b^l$ efficiently, we propose a cheap finite-difference scheme that utilizes task-specific gradients of perturbed weights. This process avoids explicit construction of the Hessian and requires only $K^2$ forward-backward passes during calibration, completing in less than 5 seconds on a single GPU.

We summarize our key contributions as follows:
\begin{enumerate}
    \item \textbf{Theoretical Rigor:} We establish a formal mathematical foundation for multi-task parameter fusion, formulating joint loss minimization as a quadratic optimization problem over task-specific local Hessian geometries.
    \item \textbf{Low-Dimensional Subspace Projection:} We prove that restricting parameter search to the $K$-dimensional subspace of task vectors allows us to model full, non-diagonal cross-parameter Hessian interactions with zero diagonal approximation and minimal computational overhead.
    \item \textbf{Training-Free Closed-Form Solution:} We derive an exact analytical closed-form solution under Ridge (Tikhonov) regularization, eliminating test-time training, transductive overfitting, and sacrificial task bias.
    \item \textbf{Empirical Characterization:} We evaluate ACM via extensive simulation sweeps (30 seeds) and physical benchmarks on Vision Transformers (ViT-Tiny) across 4 classification tasks (MNIST, FashionMNIST, CIFAR-10, SVHN). In coupled non-convex simulations, ACM outperforms PolyMerge by \textbf{+1.69\%} and AdaMerging by \textbf{+8.11\%}. On physical benchmarks, our proposed leakage-free Global-Normalized ACM (ACM-GlobalNorm) achieves a Joint Average accuracy of \textbf{57.76\%}, significantly outperforming diagonal Fisher Merging (\textbf{56.03\%}) and polynomial test-time adaptation (\textbf{38.96\%}) without any test-time optimization. We provide a transparent characterization of when local second-order approximations are optimal versus when simple uniform interpolation (Task Arithmetic at 60.72\%) exhibits regularizing advantages on fully converged, highly non-convex physical neural manifolds.
\end{enumerate}
